\documentclass[a4paper,12pt]{article}
\usepackage{fontspec}
\setmainfont{Noto Serif CJK JP}[Scale=1.0]
\usepackage{amsmath,amssymb,amsthm}
\usepackage{tikz-cd}
\usepackage{tikz}
\usetikzlibrary{positioning,shapes,arrows,decorations.pathmorphing,calc}
\usepackage{hyperref}
\usepackage[margin=25mm]{geometry}
\usepackage{listings}
\usepackage{xcolor}

\theoremstyle{definition}
\newtheorem{definition}{定義}[section]
\newtheorem{theorem}[definition]{定理}
\newtheorem{lemma}[definition]{補題}
\newtheorem{proposition}[definition]{命題}
\newtheorem{example}[definition]{例}
\newtheorem{remark}[definition]{注意}
\newtheorem{corollary}[definition]{系}

\lstdefinestyle{leanstyle}{
  basicstyle=\ttfamily\small,
  keywordstyle=\color{blue}\bfseries,
  commentstyle=\color{green!60!black},
  stringstyle=\color{red},
  showstringspaces=false,
  breaklines=true,
  frame=single,
  numbers=left,
  numberstyle=\tiny\color{gray}
}

\title{構造塔の圏と忘却関手\\
\large 圏論的視点からの形式化}
\author{Claude (Anthropic) によるTeX生成\\CODEX (OpenAI) によるLean形式化}
\date{\today}

\begin{document}

\maketitle

\begin{abstract}
本稿では、最小層付き構造塔が圏(category)を形成することを示し、その上で定義される忘却関手(forgetful functor)について解説する。忘却関手は、構造塔から台集合や添字集合への「構造を忘れる」写像であり、圏論における基本的な概念である。本稿では、Lean 4の圏論ライブラリ(Mathlib.CategoryTheory)を用いた形式化を通じて、これらの概念を厳密に理解する。また、形式化における典型的なエラーパターンとその修正方法についても議論する。
\end{abstract}

\tableofcontents

\section{導入}

\subsection{前回までの復習}

これまでに、以下の概念を学んできた:
\begin{itemize}
\item 構造塔(Structure Tower): 順序に従って単調増大する集合の族
\item 最小層付き構造塔: 各要素が初めて現れる層を明示する構造
\item 構造塔の射(morphism): 構造を保つ写像
\end{itemize}

\subsection{本稿の目的}

本稿では、これらの概念を圏論的な枠組みで統一し、以下を達成する:
\begin{enumerate}
\item 構造塔全体が圏をなすことの厳密な証明
\item 忘却関手の定義と性質の解明
\item Lean 4による形式化の実践
\item 形式化における注意点とエラー対策
\end{enumerate}

\section{圏論の基礎}

本節では、圏論の基本概念を復習する。

\begin{definition}[圏]\label{def:category}
\textbf{圏}(category)$\mathcal{C}$は以下のデータからなる:
\begin{enumerate}
\item 対象の集まり: $\text{Ob}(\mathcal{C})$
\item 射の集まり: 各対象$A, B$に対して、射の集合$\text{Hom}_{\mathcal{C}}(A, B)$
\item 合成写像: 
\[
\circ : \text{Hom}_{\mathcal{C}}(B, C) \times \text{Hom}_{\mathcal{C}}(A, B) \to \text{Hom}_{\mathcal{C}}(A, C)
\]
\item 恒等射: 各対象$A$に対して、$\text{id}_A \in \text{Hom}_{\mathcal{C}}(A, A)$
\end{enumerate}
これらは以下の公理を満たす:
\begin{itemize}
\item \textbf{結合律}: $(h \circ g) \circ f = h \circ (g \circ f)$
\item \textbf{左単位律}: $\text{id}_B \circ f = f$
\item \textbf{右単位律}: $f \circ \text{id}_A = f$
\end{itemize}
\end{definition}

\begin{example}[基本的な圏]
以下は圏の典型例である:
\begin{itemize}
\item $\mathbf{Set}$: 集合を対象、写像を射とする圏
\item $\mathbf{Grp}$: 群を対象、群準同型を射とする圏
\item $\mathbf{Top}$: 位相空間を対象、連続写像を射とする圏
\item $\mathbf{Poset}$: 半順序集合を対象、順序保存写像を射とする圏
\end{itemize}
\end{example}

\begin{definition}[関手]\label{def:functor}
圏$\mathcal{C}$から圏$\mathcal{D}$への\textbf{関手}(functor) $F : \mathcal{C} \to \mathcal{D}$は以下のデータからなる:
\begin{enumerate}
\item 対象の写像: $F : \text{Ob}(\mathcal{C}) \to \text{Ob}(\mathcal{D})$
\item 射の写像: 各$f : A \to B$に対して、$F(f) : F(A) \to F(B)$
\end{enumerate}
これらは以下の条件を満たす:
\begin{itemize}
\item \textbf{恒等射の保存}: $F(\text{id}_A) = \text{id}_{F(A)}$
\item \textbf{合成の保存}: $F(g \circ f) = F(g) \circ F(f)$
\end{itemize}
\end{definition}

\subsection{圏の図式表示}

圏の基本的な図式を以下に示す:

\begin{center}
\begin{tikzpicture}[scale=1.5]
% 対象
\node (A) at (0,0) {$A$};
\node (B) at (2,0) {$B$};
\node (C) at (4,0) {$C$};

% 射
\draw[->] (A) to node[above] {$f$} (B);
\draw[->] (B) to node[above] {$g$} (C);
\draw[->,bend right=30] (A) to node[below] {$g \circ f$} (C);

% ラベル
\node at (2,-1) {合成の可換性};
\end{tikzpicture}
\end{center}

\begin{center}
\begin{tikzpicture}[scale=1.5]
% 対象
\node (A) at (0,0) {$A$};
\node (B) at (2,0) {$B$};

% 射
\draw[->,loop above] (A) to node[above] {$\text{id}_A$} (A);
\draw[->] (A) to node[above] {$f$} (B);
\draw[->,loop above] (B) to node[above] {$\text{id}_B$} (B);

% ラベル
\node at (1,-1) {恒等射の単位律};
\end{tikzpicture}
\end{center}

\section{構造塔の圏}

\subsection{圏構造の定義}

前回定義した最小層付き構造塔と射は、実は圏を形成する。

\begin{theorem}[構造塔の圏]\label{thm:category_structure}
最小層付き構造塔全体$\mathbf{StructTowerMin}$は、以下のデータで圏をなす:
\begin{itemize}
\item 対象: 最小層付き構造塔
\item 射: $\text{Hom}(T, T') := \text{StructureTowerWithMin.Hom}(T, T')$
\item 合成: $\text{Hom.comp}$
\item 恒等射: $\text{Hom.id}$
\end{itemize}
\end{theorem}

\begin{proof}
3つの圏の公理を確認する必要がある。

\paragraph{(1) 結合律}
射$f : T \to T'$、$g : T' \to T''$、$h : T'' \to T'''$に対して、
\begin{align*}
((h \circ g) \circ f).\text{map}(x) &= (h \circ g).\text{map}(f.\text{map}(x)) \\
&= h.\text{map}(g.\text{map}(f.\text{map}(x))) \\
&= h.\text{map}((g \circ f).\text{map}(x)) \\
&= (h \circ (g \circ f)).\text{map}(x)
\end{align*}
同様に、$\text{indexMap}$についても結合律が成立する。射の外延性より、$(h \circ g) \circ f = h \circ (g \circ f)$。

\paragraph{(2) 左単位律}
射$f : T \to T'$に対して、
\begin{align*}
(\text{id}_{T'} \circ f).\text{map}(x) &= \text{id}_{T'}.\text{map}(f.\text{map}(x)) \\
&= f.\text{map}(x)
\end{align*}
したがって、$\text{id}_{T'} \circ f = f$。

\paragraph{(3) 右単位律}
同様に、$f \circ \text{id}_T = f$が成立する。
\end{proof}

\subsection{Lean 4での実装}

Lean 4では、Mathlibの圏論ライブラリを用いて以下のように実装する:

\begin{lstlisting}[style=leanstyle, language=ML]
instance : CategoryTheory.Category StructureTowerWithMin.{u, v} where
  Hom := Hom
  id := Hom.id
  comp := fun f g => Hom.comp g f  -- 注意: 順序が逆
  id_comp := by intros; rfl
  comp_id := by intros; rfl
  assoc := by intros; rfl
\end{lstlisting}

\begin{remark}[合成の順序]
Lean 4の圏論ライブラリでは、合成$\text{comp}$は「図式順」(diagrammatic order)で定義される。つまり、$f : A \to B$と$g : B \to C$に対して、$g \circ f$は$\text{comp } f \, g$と書かれる。一方、数学的な記法では通常$g \circ f$と書く。この違いに注意が必要である。
\end{remark}

\section{忘却関手}

\subsection{忘却関手の概念}

\begin{definition}[忘却関手の直観]
\textbf{忘却関手}(forgetful functor)とは、数学的構造から「構造を忘れて」より単純な対象へ写す関手である。例えば:
\begin{itemize}
\item 群から集合への忘却関手: 群の演算を忘れて台集合だけを取り出す
\item 位相空間から集合への忘却関手: 位相構造を忘れる
\item ベクトル空間から加法群への忘却関手: スカラー倍を忘れる
\end{itemize}
\end{definition}

\subsection{台集合への忘却関手}

構造塔から台集合への忘却関手を定義する。

\begin{definition}[台集合への忘却関手]\label{def:forget_carrier}
関手$U_{\text{carrier}} : \mathbf{StructTowerMin} \to \mathbf{Type}_u$を以下で定義する:
\begin{itemize}
\item 対象: $U_{\text{carrier}}(T) := T.\text{carrier}$
\item 射: $U_{\text{carrier}}(f) := f.\text{map}$ (ただし$f : T \to T'$)
\end{itemize}
\end{definition}

\begin{proposition}
$U_{\text{carrier}}$は関手の公理を満たす。
\end{proposition}

\begin{proof}
\paragraph{(1) 恒等射の保存}
$T$を任意の構造塔とする。このとき、
\begin{align*}
U_{\text{carrier}}(\text{id}_T) &= (\text{id}_T).\text{map} \\
&= \text{id}_{T.\text{carrier}} \quad \text{(定義より)}
\end{align*}

\paragraph{(2) 合成の保存}
射$f : T \to T'$、$g : T' \to T''$に対して、
\begin{align*}
U_{\text{carrier}}(g \circ f) &= (g \circ f).\text{map} \\
&= g.\text{map} \circ f.\text{map} \quad \text{(定義より)} \\
&= U_{\text{carrier}}(g) \circ U_{\text{carrier}}(f)
\end{align*}
\end{proof}

\subsection{添字集合への忘却関手}

同様に、添字集合への忘却関手も定義できる。

\begin{definition}[添字集合への忘却関手]\label{def:forget_index}
関手$U_{\text{Index}} : \mathbf{StructTowerMin} \to \mathbf{Type}_v$を以下で定義する:
\begin{itemize}
\item 対象: $U_{\text{Index}}(T) := T.\text{Index}$
\item 射: $U_{\text{Index}}(f) := f.\text{indexMap}$
\end{itemize}
\end{definition}

\begin{proposition}
$U_{\text{Index}}$は関手の公理を満たす。
\end{proposition}

\begin{proof}
台集合の場合と同様に、恒等射の保存と合成の保存を確認できる。
\end{proof}

\subsection{忘却関手の図式}

忘却関手の働きを図式で表すと以下のようになる:

\begin{center}
\begin{tikzcd}[row sep=huge, column sep=huge]
\mathbf{StructTowerMin} \arrow[r, "U_{\text{carrier}}"] \arrow[d, "U_{\text{Index}}"'] 
& \mathbf{Type}_u \\
\mathbf{Type}_v &
\end{tikzcd}
\end{center}

具体的な射の図式:

\begin{center}
\begin{tikzcd}[row sep=large, column sep=large]
T \arrow[r, "f"] \arrow[d, "U_{\text{carrier}}"'] 
& T' \arrow[d, "U_{\text{carrier}}"] \\
T.\text{carrier} \arrow[r, "f.\text{map}"'] 
& T'.\text{carrier}
\end{tikzcd}
\end{center}

この図式の可換性は、$U_{\text{carrier}}(f) = f.\text{map}$という定義そのものである。

\section{Lean 4での形式化}

\subsection{正しい実装}

Lean 4での忘却関手の正しい実装は以下のようになる:

\begin{lstlisting}[style=leanstyle, language=ML]
open CategoryTheory

def forgetCarrier : StructureTowerWithMin.{u, v} ⥤ Type u where
  obj := fun T => T.carrier
  map := fun f => f.map
  map_id := by intro T; funext x; rfl
  map_comp := by intro T T' T'' f g; funext x; rfl

def forgetIndex : StructureTowerWithMin.{u, v} ⥤ Type v where
  obj := fun T => T.Index
  map := fun f => f.indexMap
  map_id := by intro T; funext i; rfl
  map_comp := by intro T T' T'' f g; funext i; rfl
\end{lstlisting}

\subsection{形式化における注意点}

\paragraph{記法の理解}
\begin{itemize}
\item $\mathcal{C} \Rightarrow \mathcal{D}$: 圏$\mathcal{C}$から$\mathcal{D}$への関手の型
\item Leanでは$\mathcal{C} \Rightarrow \mathcal{D}$は\texttt{C ⥤ D}と書く(Unicode記号)
\item \texttt{funext}: 関数の外延性(function extensionality)を用いた証明タクティク
\end{itemize}

\paragraph{証明の構造}
関手の証明では、以下の2つの性質を示す必要がある:
\begin{enumerate}
\item \texttt{map\_id}: 恒等射の保存
\item \texttt{map\_comp}: 合成の保存
\end{enumerate}

これらの証明は、多くの場合\texttt{funext}(関数外延性)と\texttt{rfl}(反射性)で完結する。

\subsection{典型的なエラーパターン}

形式化において遭遇する典型的なエラーを以下に示す:

\paragraph{エラー1: 型の不一致}
\begin{lstlisting}[style=leanstyle, language=ML]
-- 誤り: objの型が間違っている
def forgetCarrierWrong : StructureTowerWithMin ⥤ Set where
  obj := fun T => T.carrier  -- エラー: Setではなく Type u
  ...
\end{lstlisting}

\textbf{原因}: \texttt{Type u}と\texttt{Set}は異なる型である。

\paragraph{エラー2: 合成の順序}
\begin{lstlisting}[style=leanstyle, language=ML]
-- 誤り: 合成の順序が逆
map_comp := by
  intro T T' T'' f g
  -- f : T -> T', g : T' -> T''
  -- Leanでは comp f g = g ∘ f だが、
  -- 数学的記法では g ∘ f
  funext x
  rfl  -- これは正しい
\end{lstlisting}

\textbf{注意}: Lean 4の圏論ライブラリでは、\texttt{comp f g}は数学的には$g \circ f$を意味する。

\paragraph{エラー3: 証明の不完全性}
\begin{lstlisting}[style=leanstyle, language=ML]
-- 誤り: 証明が不完全
map_id := by
  intro T
  -- funext を忘れている
  rfl  -- エラー: 関数の等式を示せない
\end{lstlisting}

\textbf{修正}: \texttt{funext}を使って関数の外延的等式に変換する必要がある。

\section{忘却関手の性質}

\subsection{忠実性と充満性}

\begin{definition}[忠実関手と充満関手]
関手$F : \mathcal{C} \to \mathcal{D}$について:
\begin{itemize}
\item $F$が\textbf{忠実}(faithful)であるとは、任意の$A, B \in \text{Ob}(\mathcal{C})$に対して、
\[
F : \text{Hom}_{\mathcal{C}}(A, B) \to \text{Hom}_{\mathcal{D}}(F(A), F(B))
\]
が単射であることをいう。
\item $F$が\textbf{充満}(full)であるとは、上の写像が全射であることをいう。
\end{itemize}
\end{definition}

\begin{proposition}
忘却関手$U_{\text{carrier}} : \mathbf{StructTowerMin} \to \mathbf{Type}_u$は忠実であるが、充満ではない。
\end{proposition}

\begin{proof}[証明のスケッチ]
\paragraph{忠実性}
2つの射$f, g : T \to T'$に対して、$U_{\text{carrier}}(f) = U_{\text{carrier}}(g)$、すなわち$f.\text{map} = g.\text{map}$であるとする。しかし、これだけでは$f = g$は言えない。実際、$f.\text{indexMap} \neq g.\text{indexMap}$である可能性がある。

したがって、$U_{\text{carrier}}$は一般には忠実ではない。

\paragraph{充満性}
任意の写像$h : T.\text{carrier} \to T'.\text{carrier}$が、ある射$f : T \to T'$の台写像$f.\text{map}$として実現されるわけではない。なぜなら、射であるためには層保存性や最小層保存性などの条件を満たす必要があるからである。
\end{proof}

\begin{remark}
上の命題の「忠実性」の主張は誤りである。実際、構造塔の射は台写像と添字写像の両方で決まるため、台写像だけでは射を一意に決定できない。正確には、$U_{\text{carrier}}$は忠実ではない。
\end{remark}

\subsection{積関手}

台写像と添字写像を同時に取り出す関手を考えることができる。

\begin{definition}[積忘却関手]
関手$U_{\times} : \mathbf{StructTowerMin} \to \mathbf{Type}_u \times \mathbf{Type}_v$を次で定義する:
\begin{itemize}
\item 対象: $U_{\times}(T) := (T.\text{carrier}, T.\text{Index})$
\item 射: $U_{\times}(f) := (f.\text{map}, f.\text{indexMap})$
\end{itemize}
\end{definition}

\begin{proposition}
$U_{\times}$は忠実関手である。
\end{proposition}

\begin{proof}
射$f, g : T \to T'$に対して、$U_{\times}(f) = U_{\times}(g)$とする。これは、
\[
(f.\text{map}, f.\text{indexMap}) = (g.\text{map}, g.\text{indexMap})
\]
を意味する。したがって、$f.\text{map} = g.\text{map}$かつ$f.\text{indexMap} = g.\text{indexMap}$である。射の外延性定理より、$f = g$が従う。
\end{proof}

\section{関手の図式による理解}

\subsection{関手の可換図式}

関手$F : \mathcal{C} \to \mathcal{D}$は、以下の可換図式で特徴付けられる:

\begin{center}
\begin{tikzcd}[row sep=huge, column sep=huge]
A \arrow[r, "f"] \arrow[d, "F"'] & B \arrow[d, "F"] \\
F(A) \arrow[r, "F(f)"'] & F(B)
\end{tikzcd}
\end{center}

\subsection{忘却関手の階層}

複数の忘却関手の関係を図式で表すと以下のようになる:

\begin{center}
\begin{tikzcd}[row sep=large, column sep=large]
& \mathbf{StructTowerMin} \arrow[dl, "U_{\text{carrier}}"'] \arrow[dr, "U_{\text{Index}}"] \arrow[d, "U_{\times}"] & \\
\mathbf{Type}_u & \mathbf{Type}_u \times \mathbf{Type}_v \arrow[l] \arrow[r] & \mathbf{Type}_v
\end{tikzcd}
\end{center}

\subsection{自然変換への展望}

2つの関手$F, G : \mathcal{C} \to \mathcal{D}$の間には、\textbf{自然変換}(natural transformation)と呼ばれる射を定義できる。これは「関手の間の射」であり、圏論のより高次の構造を形成する。

例えば、$U_{\text{carrier}}$と$U_{\text{Index}}$は異なる行き先を持つため直接比較できないが、適切な設定下で関連付けることができる。

\section{具体例と応用}

\begin{example}[群から集合への忘却関手]
群の圏$\mathbf{Grp}$から集合の圏$\mathbf{Set}$への忘却関手$U : \mathbf{Grp} \to \mathbf{Set}$は以下で定義される:
\begin{itemize}
\item 対象: $U(G) := |G|$ (群$G$の台集合)
\item 射: $U(\varphi) := \varphi$ (群準同型$\varphi$は集合の写像でもある)
\end{itemize}
これは構造塔の忘却関手と同じパターンである。
\end{example}

\begin{example}[位相空間から集合への忘却関手]
位相空間の圏$\mathbf{Top}$から$\mathbf{Set}$への忘却関手:
\begin{itemize}
\item 対象: 位相空間$(X, \tau)$を台集合$X$に写す
\item 射: 連続写像をそのまま写像として扱う
\end{itemize}
\end{example}

\begin{example}[ベクトル空間から加法群への忘却関手]
ベクトル空間の圏$\mathbf{Vect}_K$(体$K$上)から加法群の圏$\mathbf{Ab}$への忘却関手:
\begin{itemize}
\item 対象: ベクトル空間$V$をその加法群構造$(V, +)$に写す
\item 射: 線形写像を群準同型として扱う
\end{itemize}
\end{example}

\section{発展的話題}

\subsection{随伴関手}

忘却関手には、しばしば「自由関手」(free functor)と呼ばれる左随伴(left adjoint)が存在する。

\begin{definition}[随伴]\label{def:adjunction}
関手$F : \mathcal{C} \to \mathcal{D}$と$U : \mathcal{D} \to \mathcal{C}$について、$F$が$U$の\textbf{左随伴}であるとは、自然な同型
\[
\text{Hom}_{\mathcal{D}}(F(C), D) \cong \text{Hom}_{\mathcal{C}}(C, U(D))
\]
が存在することをいう。記号で$F \dashv U$と書く。
\end{definition}

\begin{example}[自由群]
集合の圏$\mathbf{Set}$から群の圏$\mathbf{Grp}$への自由群関手$F$は、忘却関手$U : \mathbf{Grp} \to \mathbf{Set}$の左随伴である:
\[
F \dashv U
\]
これは、集合$X$と群$G$に対して、
\[
\text{Hom}_{\mathbf{Grp}}(F(X), G) \cong \text{Hom}_{\mathbf{Set}}(X, U(G))
\]
が成立することを意味する。
\end{example}

\subsection{極限と余極限}

忘却関手は、しばしば極限や余極限を保存する。

\begin{definition}[極限の保存]
関手$F : \mathcal{C} \to \mathcal{D}$が\textbf{極限を保存する}とは、$\mathcal{C}$における任意の極限$\varprojlim D_i$に対して、
\[
F(\varprojlim D_i) \cong \varprojlim F(D_i)
\]
が成立することをいう。
\end{definition}

多くの忘却関手は極限を保存するが、余極限は保存しないことが多い。

\subsection{表現可能関手}

構造塔の圏において、特定の対象$T_0$を固定したとき、関手
\[
\text{Hom}(T_0, -) : \mathbf{StructTowerMin} \to \mathbf{Set}
\]
は\textbf{表現可能関手}(representable functor)と呼ばれる。これは圏論における重要な概念である。

\section{形式化の実践的側面}

\subsection{証明の自動化}

Lean 4では、多くの圏論的証明が自動化されている:

\begin{lstlisting}[style=leanstyle, language=ML]
-- 結合律の証明
assoc := by intros; rfl

-- 単位律の証明
id_comp := by intros; rfl
comp_id := by intros; rfl
\end{lstlisting}

これらの証明が\texttt{rfl}(反射性)だけで完結するのは、定義を展開すると自動的に等式が成立するためである。

\subsection{型クラスの活用}

Lean 4では、型クラス(type class)を用いて圏構造を定義する:

\begin{lstlisting}[style=leanstyle, language=ML]
instance : CategoryTheory.Category StructureTowerWithMin where
  ...
\end{lstlisting}

この\texttt{instance}宣言により、Leanは構造塔が圏であることを自動的に認識し、圏論の一般的な定理を適用できるようになる。

\subsection{宇宙レベルの管理}

形式化では、宇宙レベル(universe level)の管理が重要である:

\begin{lstlisting}[style=leanstyle, language=ML]
universe u v

structure StructureTowerWithMin where
  carrier : Type u
  Index : Type v
  ...
\end{lstlisting}

これにより、異なる大きさの型を扱うことができ、ラッセルのパラドックスなどの集合論的矛盾を回避できる。

\section{まとめ}

本稿では、構造塔の圏論的扱いと忘却関手について解説した。主要な成果は以下の通りである:

\begin{itemize}
\item 構造塔が圏をなすことの厳密な証明
\item 台集合および添字集合への忘却関手の定義
\item Lean 4による形式化の実装
\item 形式化における典型的なエラーパターンとその対策
\item 忠実性・充満性などの関手の性質
\item 随伴関手や表現可能関手への展望
\end{itemize}

圏論は現代数学の共通言語であり、構造塔の理論を圏論的に定式化することで、より広い数学的文脈でその性質を理解できる。Lean 4による形式化は、これらの抽象的な概念を機械的に検証可能な形で記述する強力な手段である。

\section*{参考文献}

\begin{enumerate}
\item Saunders Mac Lane, \textit{Categories for the Working Mathematician}, Springer, 1971.
\item Emily Riehl, \textit{Category Theory in Context}, Dover, 2016.
\item Tom Leinster, \textit{Basic Category Theory}, Cambridge University Press, 2014.
\item The mathlib Community, \textit{The Lean Mathematical Library}, \url{https://github.com/leanprover-community/mathlib4}
\item N. Bourbaki, \textit{Elements of Mathematics: Theory of Sets}, Springer, 2004.
\item 圏論の基礎, 山口佳三, 日本評論社, 2010.
\item 圏論の歩き方, 圏論の会(編), 日本評論社, 2015.
\end{enumerate}

\appendix

\section{Lean 4コードの全文}

本稿で解説した構造塔の圏と忘却関手の形式化の全体は以下の通りである:

\begin{lstlisting}[style=leanstyle, language=ML]
import Mathlib.Data.Set.Basic
import Mathlib.Order.Basic
import Mathlib.CategoryTheory.Category.Basic
import Mathlib.CategoryTheory.Functor.Basic
import Mathlib.CategoryTheory.Types.Basic

universe u v

structure StructureTowerWithMin where
  carrier : Type u
  Index : Type v
  [indexPreorder : Preorder Index]
  layer : Index → Set carrier
  covering : ∀ x : carrier, ∃ i : Index, x ∈ layer i
  monotone : ∀ {i j : Index}, i ≤ j → layer i ⊆ layer j
  minLayer : carrier → Index
  minLayer_mem : ∀ x, x ∈ layer (minLayer x)
  minLayer_minimal : ∀ {x i}, x ∈ layer i → minLayer x ≤ i

namespace StructureTowerWithMin

instance instIndexPreorder (T : StructureTowerWithMin.{u, v}) : 
    Preorder T.Index := T.indexPreorder

variable {T T' T'' : StructureTowerWithMin.{u, v}}

structure Hom (T T' : StructureTowerWithMin.{u, v}) where
  map : T.carrier → T'.carrier
  indexMap : T.Index → T'.Index
  indexMap_mono : ∀ {i j : T.Index}, i ≤ j → 
    indexMap i ≤ indexMap j
  layer_preserving : ∀ (i : T.Index) (x : T.carrier),
    x ∈ T.layer i → map x ∈ T'.layer (indexMap i)
  minLayer_preserving : ∀ x, 
    T'.minLayer (map x) = indexMap (T.minLayer x)

def Hom.id (T : StructureTowerWithMin.{u, v}) : Hom T T where
  map := _root_.id
  indexMap := _root_.id
  indexMap_mono := fun hij => hij
  layer_preserving := fun _ _ hx => hx
  minLayer_preserving := fun _ => rfl

def Hom.comp (g : Hom T' T'') (f : Hom T T') : Hom T T'' where
  map := fun x => g.map (f.map x)
  indexMap := fun i => g.indexMap (f.indexMap i)
  indexMap_mono := fun hij => 
    g.indexMap_mono (f.indexMap_mono hij)
  layer_preserving := fun i x hx =>
    g.layer_preserving (f.indexMap i) (f.map x) 
      (f.layer_preserving i x hx)
  minLayer_preserving := by
    intro x
    calc T''.minLayer (g.map (f.map x))
        = g.indexMap (T'.minLayer (f.map x)) := by 
          rw [g.minLayer_preserving]
      _ = g.indexMap (f.indexMap (T.minLayer x)) := by 
          rw [f.minLayer_preserving]

instance : CategoryTheory.Category 
    StructureTowerWithMin.{u, v} where
  Hom := Hom
  id := Hom.id
  comp := fun f g => Hom.comp g f
  id_comp := by intros; rfl
  comp_id := by intros; rfl
  assoc := by intros; rfl

open CategoryTheory

def forgetCarrier : StructureTowerWithMin.{u, v} ⥤ Type u where
  obj := fun T => T.carrier
  map := fun f => f.map
  map_id := by intro T; funext x; rfl
  map_comp := by intro T T' T'' f g; funext x; rfl

def forgetIndex : StructureTowerWithMin.{u, v} ⥤ Type v where
  obj := fun T => T.Index
  map := fun f => f.indexMap
  map_id := by intro T; funext i; rfl
  map_comp := by intro T T' T'' f g; funext i; rfl

end StructureTowerWithMin
\end{lstlisting}

\section{演習問題}

\begin{enumerate}
\item 定理\ref{thm:category_structure}の結合律を、添字写像についても詳細に証明せよ。

\item 忘却関手$U_{\text{carrier}}$が忠実でないことを、具体例を用いて示せ。

\item 積忘却関手$U_{\times}$が関手の公理(恒等射の保存と合成の保存)を満たすことを証明せよ。

\item 構造塔の圏において、同型射(isomorphism)を特徴付けよ。すなわち、射$f : T \to T'$が同型となる条件を述べよ。

\item 群の圏$\mathbf{Grp}$から集合の圏$\mathbf{Set}$への忘却関手が、積(product)を保存することを示せ。

\item 自由群関手$F : \mathbf{Set} \to \mathbf{Grp}$が忘却関手$U : \mathbf{Grp} \to \mathbf{Set}$の左随伴であることの具体的な意味を説明せよ。
\end{enumerate}

\end{document}
